\chapter{Inside the Engine Room: Virtual Machine Architecture \& Bytecodes}

\section{What is a Virtual Machine?}

\begin{kidbox}{The Imaginary Computer Inside Your Computer}
Have you ever played Minecraft or a video game where you can build an arcade machine inside the game?
You are sitting in front of your real television, playing a game, and inside that game, your character is playing another little game!
A \textbf{Virtual Machine} is an imaginary computer built completely out of software code that lives inside your real computer.
It has its own imaginary brain, its own imaginary instructions, and its own imaginary memory!
\end{kidbox}

\section{The Pancake Stack Architecture}

\begin{kidbox}{Stack of Fluffy Pancakes}
Imagine a stack of pancakes on a plate at breakfast:
\begin{itemize}
    \item If you want to add a pancake, you put it on \textbf{TOP}. We call this \textbf{PUSH}!
    \item If you want to eat a pancake, you take it from the \textbf{TOP}. We call this \textbf{POP}!
    \item You can never pull a pancake from the middle without making the whole stack fall over!
\end{itemize}
The Shae Virtual Machine is a \textbf{Stack Machine}! 
To add $10 + 20$:
\begin{enumerate}
    \item Push 10 onto the pancake stack.
    \item Push 20 onto the pancake stack.
    \item Execute \texttt{OpAdd}: The VM pops 20, pops 10, adds them together ($10 + 20 = 30$), and pushes 30 back onto the top of the stack!
\end{enumerate}
\end{kidbox}

\begin{deepbox}{The Shae Opcode Instruction Set}
The Shae compiler transforms AST nodes into high-density 8-bit bytecode instructions (\texttt{src/chunk.rs}, \texttt{src/compiler.rs}, \texttt{src/vm.rs}):

\begin{longtable}{lll}
\toprule
\textbf{Opcode} & \textbf{Operands} & \textbf{Stack Transformation \& Description} \\
\midrule
\endhead
\texttt{OpConstant} & \texttt{idx: u16} & $[\dots] \to [\dots, \text{constants}[idx]]$ \\
\texttt{OpPop} & None & $[\dots, v] \to [\dots]$ \\
\texttt{OpDup} & None & $[\dots, v] \to [\dots, v, v]$ \\
\texttt{OpAdd} & None & $[\dots, a, b] \to [\dots, a + b]$ \\
\texttt{OpSub} & None & $[\dots, a, b] \to [\dots, a - b]$ \\
\texttt{OpMul} & None & $[\dots, a, b] \to [\dots, a \times b]$ \\
\texttt{OpDiv} & None & $[\dots, a, b] \to [\dots, a / b]$ \\
\texttt{OpModulo} & None & $[\dots, a, b] \to [\dots, a \pmod b]$ \\
\texttt{OpEqual} & None & $[\dots, a, b] \to [\dots, a == b]$ \\
\texttt{OpGreater} & None & $[\dots, a, b] \to [\dots, a > b]$ \\
\texttt{OpLess} & None & $[\dots, a, b] \to [\dots, a < b]$ \\
\texttt{OpNot} & None & $[\dots, v] \to [\dots, \neg v]$ \\
\texttt{OpBitAnd} & None & $[\dots, a, b] \to [\dots, a \ \& \ b]$ \\
\texttt{OpBitOr} & None & $[\dots, a, b] \to [\dots, a \mid b]$ \\
\texttt{OpBitXor} & None & $[\dots, a, b] \to [\dots, a \oplus b]$ \\
\texttt{OpBitShl} & None & $[\dots, a, b] \to [\dots, a \ll b]$ \\
\texttt{OpBitShr} & None & $[\dots, a, b] \to [\dots, a \gg b]$ \\
\texttt{OpGetLocal} & \texttt{slot: u16} & $[\dots] \to [\dots, \text{stack}[\text{base} + slot]]$ \\
\texttt{OpSetLocal} & \texttt{slot: u16} & $[\dots, v] \to [\dots]$ (writes to local slot) \\
\texttt{OpGetGlobal} & \texttt{name\_idx: u16} & $[\dots] \to [\dots, \text{globals}[\text{name}]]$ \\
\texttt{OpSetGlobal} & \texttt{name\_idx: u16} & $[\dots, v] \to [\dots]$ \\
\texttt{OpJump} & \texttt{offset: u16} & Unconditional jump: $ip \gets ip + offset$ \\
\texttt{OpJumpIfFalse}& \texttt{offset: u16} & Jump if top of stack is falsy \\
\texttt{OpLoop} & \texttt{offset: u16} & Backwards loop jump: $ip \gets ip - offset$ \\
\texttt{OpClosure} & \texttt{fn\_idx: u16} & Pushes closure wrapping function and upvalues \\
\texttt{OpCall} & \texttt{argc: u8} & Invokes callable with $argc$ arguments \\
\texttt{OpReturn} & None & Returns top of stack to caller frame \\
\texttt{OpArray} & \texttt{count: u16} & Pops $count$ values and creates dynamic array \\
\texttt{OpMap} & \texttt{pairs: u16} & Pops $pairs \times 2$ keys and values to form map \\
\bottomrule
\end{longtable}
\end{deepbox}

\section{Disassembling Bytecode with \texttt{shae disasm}}

You can inspect the exact compiled bytecode of any Shae program using the CLI disassembler:
\begin{lstlisting}[language=bash]
shae disasm program.shae
\end{lstlisting}

Output:
\begin{lstlisting}[language=shaelang]
== Disassembly: main ==
Offset  Line  OpCode            Operand
---------------------------------------------
0000    1     OpConstant        #0 (100)
0003    1     OpConstant        #1 (50)
0006    1     OpAdd
0007    1     OpSetLocal        slot #0 (score)
0010    2     OpGetLocal        slot #0
0013    2     OpPrint
0014    2     OpReturn
\end{lstlisting}
